% Proof-complete for the three pinned middle-column theorems.
% Earlier draft evidence, roadmap, and prior-art limits are retained as historical material.
% Original source packet: BRIEF.md, REPORT-consolidation.md, REPORT-threshold.md.
% Refreshed from the supplied frozen proof sources, PIN.md, and crossed-review record.
\documentclass[11pt]{article}
\usepackage[T1]{fontenc}
\usepackage[utf8]{inputenc}
\usepackage{lmodern}
\usepackage[margin=1in]{geometry}
\usepackage{amsmath,amssymb,amsthm}
\usepackage{booktabs,longtable,array}
\usepackage[hidelinks]{hyperref}
\hypersetup{pdftitle={Binary submasks and the middle column of Rules 126 and 129},
  pdfauthor={Derek Blair}}
\newtheorem{targettheorem}{Target Theorem}
\newcommand{\pc}{\operatorname{popcount}}
\newcommand{\sub}{\preceq}
\newcommand{\ind}{\mathbf{1}}
\newcommand{\Z}{\mathbb{Z}}
\newcommand{\N}{\mathbb{N}}
\newcommand{\seq}[1]{\text{\texttt{#1}}}
\newcommand{\missing}{\textit{not supplied}}
\setlength{\emergencystretch}{2em}
\pagestyle{myheadings}
\markright{Pinned middle-column proofs complete}

\title{Binary submasks and the middle column\protect\\
of Rules 126 and 129\thanks{Proof-complete for the three pinned middle-column theorems;
finite computational evidence and prior-art limitations retained.}}
\author{Derek Blair}
\date{28 October 2026}

\begin{document}
\maketitle

\begin{abstract}
We formulate a target characterization of the rows of elementary cellular
automaton Rule 126 from a single live cell, and its consequence for the middle
column of Rule 129 under an exact complementation reduction. The candidate
support is obtained by reflecting a union indexed by the binary submasks of
the lower part of the generation number. For a generation with $p$ set bits,
this gives $2^{p+1}$ live cells except at $n=2^p-1$, where the origin is
counted once and the candidate row is full. The resulting middle-column
identity is $\seq{A267442}(n)=\seq{A043545}(n)$ for $n\geq1$, through
$\seq{A036987}$; the equality fails at $n=0$. We report finite computational
evidence for popcounts $1$ through $18$, including the second-smallest
$p$-bit threshold, and specify a binomial-parity and Lucas-theorem roadmap.
The three pinned middle-column theorems now have completed Lean proofs;
the broader row-support target and roadmap are retained as the earlier
draft record. The sequence identity is already an OEIS conjecture, and related
row-count and binomial formulas are prior art; no novelty is claimed for
those assertions.
\end{abstract}

\noindent\textbf{Status and scope.}
The three pinned middle-column theorems now have a completed Lean proof,
as described in Section~\ref{sec:formal-proof}. This proof-complete status
concerns the Rule 126 middle column and the positive-index Rule 129
identity. The original row-support target, proof roadmap, computational
evidence, and prior-art discussion below are retained as the earlier
draft record; their references to missing proof artifacts describe that
record's original status. In particular, the broader row-support and
row-count formulation is not presented here as an additional pinned
formal theorem. Completion of the finite checks remains distinct from
a universal proof, and no novelty claim or computational coverage
claim is strengthened.

\section{Introduction}

Elementary cellular automata update a bi-infinite binary row using the same
Boolean function of three neighboring cells at every position. In the
standard numbering convention, Rule 126 returns zero on $000$ and $111$ and
one on the remaining six triples. Rule 129 has the complementary output
table. Wolfram's \emph{A New Kind of Science} describes these automata and
the relation of Rule 126 to Rule 90~\cite{wolfram126}; it also discusses the
power-of-two pattern in the center column of Rule 129~\cite{wolfram129}.

The sequence of center-cell states for Rule 129 from a single live cell is
catalogued as OEIS A267442~\cite{oeis267442}. The two comparison sequences are
\begin{align}
 r(n)&:=\seq{A036987}(n)
       =\ind_{\{n=2^j-1\text{ for some }j\geq0\}},\label{eq:rueppel}\\
 q(n)&:=\seq{A043545}(n)=1-r(n),\qquad n\geq0.
       \label{eq:q}
\end{align}
The first is the Fredholm--Rueppel sequence; the second is its binary
complement~\cite{oeis036987,oeis043545}. The intended headline identity is
\begin{equation}
 \boxed{\seq{A267442}(n)=\seq{A043545}(n)\quad(n\geq1).}
 \label{eq:headline}
\end{equation}
The positive-index qualification is essential. A267442 has value $1$ at
index zero, whereas A043545 has value $0$ there. Richard J.\ Mathar's
17 January 2016 comment on A267442 already states
\eqref{eq:headline} as a conjecture~\cite{oeis267442}.

Our purpose is to organize a precise row-support statement, the finite
evidence supplied for it, and the obligations of a subsequent formal
derivation. The middle column is the simplest consequence of the proposed
support formula. The threshold
\[
 T_p=3\cdot2^{p-1}-1
\]
has an elementary interpretation: it is the second-smallest integer with
$p$ set bits. Thus the apparent gap below the threshold contains only the
single smaller $p$-bit integer $2^p-1$, and no other candidate generations.

\section{Conventions and the exact complementation reduction}
\label{sec:conventions}

Let $B_n,S_n:\Z\longrightarrow\{0,1\}$ be the Rule 126 and Rule 129 rows,
respectively, with the \emph{same} initial condition
\begin{equation}
 B_0(x)=S_0(x)=\ind_{\{x=0\}}.
 \label{eq:seed}
\end{equation}
Writing $f_{126},f_{129}$ for the two local maps, the updates are
\[
 B_{n+1}(x)=f_{126}(B_n(x-1),B_n(x),B_n(x+1)),
\]
and likewise for $S$ with $f_{129}$. Explicitly,
\[
 f_{126}(a,b,c)=\ind_{\{(a,b,c)\notin\{(0,0,0),(1,1,1)\}\}},
 \qquad f_{129}=1-f_{126}.
\]
Both maps are invariant under complementing all three inputs. If $C$ is
pointwise complementation and $F_{126},F_{129}$ are the global update maps,
these truth-table identities give
\[
 F_{129}=C F_{126},\qquad F_{126}C=F_{126}.
\]
Starting at time one and iterating these equalities yields the exact reduction
\begin{equation}
 S_n(x)=1-B_n(x)\qquad(n\geq1,\ x\in\Z).
 \label{eq:complement}
\end{equation}
There is no alternating sign in time. At time zero, \eqref{eq:seed} applies
instead.

The Rule 126 row vanishes outside the cone $|x|\leq n$ and is symmetric
under $x\mapsto-x$. We write
\[
 C_B(n)=\#\{x\in\Z:B_n(x)=1\},\qquad
 C_S^{\mathrm{cone}}(n)=\#\{x\in[-n,n]\cap\Z:S_n(x)=1\}.
\]
For $n\geq1$, the Rule 129 background is one outside the cone, so a count
of all its live cells would be infinite. The finite convention used here is
\begin{equation}
 C_S^{\mathrm{cone}}(n)=2n+1-C_B(n)\qquad(n\geq1).
 \label{eq:conecount}
\end{equation}
This makes explicit the cone interpretation needed for the reports'
Mersenne assertion $\texttt{count\_S}=0$; their symbol is not separately
defined in the supplied packet.

\section{Completed formal proof of the pinned middle-column statements}
\label{sec:formal-proof}

The proof closure establishes the three pinned declarations in
\texttt{Mathproofs.Rule129MiddleColumn}:
\texttt{rule126\_middle\_column\_iff},
\texttt{rule126\_middle\_column\_eq\_a036987}, and
\texttt{a267442\_eq\_a043545}.
The first characterizes the Rule 126 center cell at every natural
generation by the Mersenne condition. The second identifies that
middle column with A036987 at every natural generation. The third
establishes the stated equality of A267442 and A043545 at positive
indices, with the original seed, dynamics, definitions, and domain.

The module structure consists of \texttt{Rule129ClosedForm},
\texttt{Rule129Lucas}, and \texttt{Rule129Digits}, followed by the
\texttt{Rule129MiddleColumn} closure. ClosedForm verifies the seed and
Rule 126 recurrence for the binomial-parity expression. Lucas supplies
the central-binomial parity criterion; Digits supplies the binary-digit
characterization. MiddleColumn links the closed form to the pinned
orbit and combines these ingredients with the positive-time
complementation reduction. The axiom allowlist is
\texttt{[propext, Classical.choice, Quot.sound]}.

The supplied crossed-review history is M1 v2 PASS, M2 PASS, M3 PASS,
and M4 PASS, with reviewer pool \texttt{chatgpt-proton} and author
pool \texttt{chatgpt-live}. The supplied \texttt{REVIEW-M4.md} records
the independent M4 review and treats the earlier helper PASS
statuses as supplied context; it does not claim to reproduce those
earlier reviews. Verification performed during this PR assembly
belongs to the author lane and is not an additional independent
review.

\section{The target row characterization and the middle column}
\label{sec:target}

For a positive integer $n$, define
\[
 p=\pc(n),\qquad H(n)=2^{\lfloor\log_2 n\rfloor},\qquad L(n)=n-H(n).
\]
Here $\pc$ counts the one-bits in the binary expansion, $H$ is the leading
power of two, and $L$ is the \emph{lower part}. For nonnegative integers,
$k\sub L$ means that every set bit of $k$ is a set bit of $L$; equivalently,
the bitwise intersection of $k$ and $L$ is $k$. The lower part has $p-1$
set bits and hence $2^{p-1}$ submasks. Define
\begin{equation}
 U(n)=\bigcup_{k\sub L(n)}\{n-2k,\ n-2k-1\}.
 \label{eq:union}
\end{equation}
This is a set of candidate \emph{absolute positions}; reflection is required
to obtain positions in the whole row.

\begin{targettheorem}[Lower-part support and Mersenne exception;
pending proof]\label{thm:target}
Let $p\geq1$ and $n\geq1$ be integers with $\pc(n)=p$. If
\begin{equation}
 n\geq3\cdot2^{p-1}-1,\qquad n\ne2^p-1,
 \label{eq:hyp}
\end{equation}
then the proposed generation-$n$ Rule 126 row is characterized by
\begin{equation}
 B_n(x)=1\quad\Longleftrightarrow\quad |x|\in U(n)
 \qquad(x\in\Z).
 \label{eq:support}
\end{equation}
The union has $2^p$ distinct positive values, and
\begin{equation}
 C_B(n)=2^{p+1},\qquad B_n(0)=0,\qquad S_n(0)=1.
 \label{eq:ordinarycount}
\end{equation}
At the exceptional generation $n=2^p-1$, the same union description gives
\begin{equation}
 U(n)=\{0,1,\ldots,n\},\qquad
 B_n(x)=\ind_{\{|x|\leq n\}},\qquad C_B(n)=2^{p+1}-1.
 \label{eq:mersenne}
\end{equation}
The value zero is produced exactly once in \eqref{eq:union}, from
$k=L(n)$ in the term $n-2k-1$. Consequently $B_n(0)=1$, $S_n(0)=0$,
and $C_S^{\mathrm{cone}}(n)=0$.
\end{targettheorem}

\noindent\textbf{Logical status.}
Target Theorem~\ref{thm:target} states the intended universal result, for
arbitrary $p\geq1$. It is not established here by the reported finite
checks. Those checks concern specified samples in the tiers $1\leq p\leq18$;
they do not quantify over every generation of even those tiers. The
redundant exclusion in \eqref{eq:hyp} is retained to match the consolidated
statement supplied for this draft.

\subsection{Why the threshold is the second \texorpdfstring{$p$}{p}-bit integer}

The smallest integer with $p$ set bits is $M_p=2^p-1$, using positions
$0,\ldots,p-1$. Any other integer with $p$ set bits must use a position
at least $p$. Its other $p-1$ bits contribute at least $2^{p-1}-1$.
Equality is attained by using positions $0,\ldots,p-2,p$, giving
\[
 T_p=2^p+(2^{p-1}-1)=3\cdot2^{p-1}-1.
\]
For $p=1$ the lower list of positions is empty, so the same reasoning gives
$M_1=1$ and $T_1=2$. In particular, among integers of popcount $p$,
\begin{equation}
 n\geq T_p\quad\Longleftrightarrow\quad n\ne M_p.
 \label{eq:second}
\end{equation}
This is an ordering argument about binary integers, independent of any
cellular-automaton assertion. There is no interval of additional $p$-bit
test cases between $M_p$ and $T_p$.

\subsection{The conditional middle-column consequence}

Once Target Theorem~\ref{thm:target}, including its exceptional case, is
established, its center-cell consequence is
\begin{equation}
 B_n(0)=r(n)\quad(n\geq0),\qquad
 S_n(0)=q(n)\quad(n\geq1).
 \label{eq:center}
\end{equation}
The first equality at $n=0$ follows separately from the seed. Equations
\eqref{eq:complement} and \eqref{eq:rueppel} then give the intended
positive-index identity \eqref{eq:headline}. The endpoint distinction is
visible in the sequence values:
\begin{center}
\begin{tabular}{lrrrrrrrrr}
\toprule
$n$ & 0&1&2&3&4&5&6&7&8\\
\midrule
A036987 & 1&1&0&1&0&0&0&1&0\\
A043545 & 0&0&1&0&1&1&1&0&1\\
A267442 & 1&0&1&0&1&1&1&0&1\\
\bottomrule
\end{tabular}
\end{center}
These initial values illustrate the indexing; a finite matching prefix
does not settle \eqref{eq:headline}.

\section{Proof roadmap: obligations and established ingredients}
\label{sec:roadmap}

This section is a roadmap, not a proof of
Target Theorem~\ref{thm:target}. The distinction is between the unresolved
connection to the nonlinear evolution and the subsequent binary algebra.

\paragraph{Step 1: fix the automata and the reduction.}
Formalize \eqref{eq:seed}, the two truth tables, finite propagation, reflection
symmetry, and \eqref{eq:complement}. The Boolean reduction in
Section~\ref{sec:conventions} is elementary; its Lean implementation and
the positive-time domain still have to be linked to the eventual statement.

\paragraph{Step 2: connect the evolution to binomial parity (pending).}
A candidate closed form consistent with the lower-part union is
\begin{equation}
 \widehat B_n(x)=
 \begin{cases}
 \displaystyle\binom{n}{\left\lfloor(n-|x|)/2\right\rfloor}\bmod2,
       & |x|\leq n,\\[6pt]
 0,    & |x|>n.
 \end{cases}
 \label{eq:binomial}
\end{equation}
The upper binomial index is $n$. The key missing obligation is
$B_n(x)=\widehat B_n(x)$ for all $n$ and $x$, including the cone edges,
the center, both spatial signs, and the seed. One possible route uses
Rule 90 or a suitable auxiliary recurrence, followed by Pascal's identity
modulo two. Rule 126 itself is nonlinear, so its update cannot simply be
replaced by XOR. This candidate is an explicit roadmap formulation; the
two supplied reports do not contain a frozen statement or a derivation of
\eqref{eq:binomial}. A related binomial formula for the right half of the
Rule 126 triangle is already catalogued as A127872~\cite{oeis127872}.

\paragraph{Step 3: apply the established Lucas criterion.}
Lucas's theorem modulo a prime, here specialized to the prime $2$, gives
\begin{equation}
 \binom{n}{k}\equiv\prod_{i\geq0}\binom{n_i}{k_i}\pmod2,
 \qquad
 \binom{n}{k}\text{ odd}\quad\Longleftrightarrow\quad k\sub n,
 \label{eq:lucas}
\end{equation}
where $n_i,k_i\in\{0,1\}$ are binary digits~\cite{lucas}. For $n\geq1$
and $|x|\leq n$, set $k=\lfloor(n-|x|)/2\rfloor$. Then
\[
 0\leq k\leq\lfloor n/2\rfloor<H(n).
\]
Thus $k\sub n$ is equivalent to $k\sub L(n)$. The floor equation is in
turn equivalent to $|x|=n-2k$ or $|x|=n-2k-1$. Accordingly Lucas's theorem
converts \eqref{eq:binomial} into the union \eqref{eq:union}. This
conversion does not discharge Step 2.

\paragraph{Step 4: count the candidate union.}
Within either family $n-2k$ or $n-2k-1$, different submasks give different
values; between the two families, parity differs. Hence $U(n)$ has $2^p$
distinct values. Its minimum is
\begin{equation}
 \min U(n)=n-2L(n)-1=H(n)-L(n)-1\geq0.
 \label{eq:min}
\end{equation}
It is zero exactly when $L(n)=H(n)-1$, or $n$ is Mersenne. Every positive
absolute position yields two reflected cells, while zero yields one. These
set-theoretic facts explain the two candidate counts and the center
criterion. They are facts about $U(n)$; they apply to the automaton after
the missing support equality is supplied.

\paragraph{Step 5: complete the formal bridge.}
The Lean development must identify its evolution, generation numbering,
binary-popcount and submask definitions, integer coordinates, and natural
binomial indices with the conventions above. It must then derive the
exceptional case, the all-popcount assertion, and the positive-index OEIS
identity from the actual row theorem. No Lean source, theorem identifier,
dependency audit, or successful Lean build is part of the present evidence
packet. The numbered target and equations here are editorial references,
not claims that corresponding Lean declarations already exist.

\section{Computational evidence and its exact limits}
\label{sec:evidence}

\subsection{Provenance of the supplied evidence}

The two supplied reports are dated 27 October 2026. The consolidation
report~\cite{report-consolidation} records a review of tranche output files,
not a new calculation. For the power-of-two tier and popcounts $2$ through
$18$, it records the terminal markers \texttt{ALL CHECKS PASS} and
\texttt{TOTAL failures: 0}. The popcount-$18$ output has both markers twice:
once for sections (A)--(E), and once for the (F) bitset run. The supporting
\texttt{mersenne-rows}, \texttt{lucas-submask}, \texttt{diagonals}, and
\texttt{rowcounts} tranches also carry \texttt{ALL CHECKS PASS}. The
\texttt{a1a4-truthtables}, \texttt{a3lucas}, \texttt{closedform}, and
\texttt{recurrence} tranches use different reporting conventions and are
not cited there as tier evidence.

The threshold report~\cite{report-threshold} describes a separate
\texttt{compute.py} run using Python and NumPy, with exit status zero and
the same success markers. It states that no Lean was invoked. The present
draft uses these two reports as evidence summaries; their underlying
scripts and raw tranche logs were not supplied with this writing task.

\subsection{All 18 tiers, without extrapolated ranges}

For every tier $1\leq p\leq18$, the consolidation reports agreement with
\eqref{eq:union} and the non-Mersenne count on the generations checked in
that tranche, as well as the exceptional full row. It also reports the
converse test
\begin{equation}
 C_B(n)=2^{p+1}
 \quad\Longleftrightarrow\quad
 \bigl(\pc(n)=p\ \text{and}\ n\ne2^p-1\bigr)
 \label{eq:converse}
\end{equation}
on each tranche's recorded converse interval. Only the largest such
interval is reproduced: $p=18$, $2\leq n\leq512000$. The other intervals
are not specified. In particular, the upper endpoint $512000$ must not be
applied to every tier.

Table~\ref{tab:tiers} reproduces every threshold and explicitly separates
these unknown ranges. The count column evaluates the reported formula
$2^{p+1}$; it is not a sample count. ``Not supplied'' means that the two
input reports omit the relevant range, not that the original tranche
performed no such check.

\begingroup
\small
\setlength{\tabcolsep}{4pt}
\begin{longtable}{@{}rrr >{\raggedright\arraybackslash}p{2.35cm}
 >{\raggedright\arraybackslash}p{3.05cm}
 >{\raggedright\arraybackslash}p{2.15cm}@{}}
\caption{Reported popcount tiers and bounds present in the source packet.}
\label{tab:tiers}\\
\toprule
$p$ & $T_p$ & $2^{p+1}$ & Analytic row sample & Converse interval & Simulation range\\
\midrule
\endfirsthead
\toprule
$p$ & $T_p$ & $2^{p+1}$ & Analytic row sample & Converse interval & Simulation range\\
\midrule
\endhead
\bottomrule
\endfoot
1 & 2 & 4 & \missing & \missing & \missing\\
2 & 5 & 8 & \missing & \missing & \missing\\
3 & 11 & 16 & \missing & \missing & \missing\\
4 & 23 & 32 & \missing & \missing & \missing\\
5 & 47 & 64 & \missing & \missing & \missing\\
6 & 95 & 128 & \missing & \missing & \missing\\
7 & 191 & 256 & \missing & \missing & \missing\\
8 & 383 & 512 & \missing & \missing & \missing\\
9 & 767 & 1024 & \missing & \missing & \missing\\
10 & 1535 & 2048 & \missing & \missing & \missing\\
11 & 3071 & 4096 & \missing & \missing & \missing\\
12 & 6143 & 8192 & \missing & \missing & \missing\\
13 & 12287 & 16384 & \missing & \missing & \missing\\
14 & 24575 & 32768 & \missing & \missing & \missing\\
15 & 49151 & 65536 & \missing & \missing & \missing\\
16 & 98303 & 131072 & \missing & \missing & \missing\\
17 & 196607 & 262144 & \missing & \missing & \missing\\
18 & 393215 & 524288 & 19 values; see below & $2\leq n\leq512000$ & \missing\\
\end{longtable}
\endgroup

The threshold report identifies the $p=18$ generation list with the
popcount-$18$ tranche's list: the 19 integers having 18 set bits among
positions $0,\ldots,18$. Equivalently this explicitly specified sample is
\begin{equation}
 E_{18}=\{524287-2^j:0\leq j\leq18\}.
 \label{eq:e18}
\end{equation}
The report describes these as 19 values $\leq524288$. This is a sparse
list, not every generation up to that bound; its exact endpoints, from
\eqref{eq:e18}, are $262143$ and $524286$. The supplied packet does not
identify the algorithm or coverage of every individual test on that list.
In particular it does not give a bitset-simulation interval, even though
it records completion of a bitset section. No analogous analytic row
lists for the other 17 main tranches are reproduced.

\subsection{Threshold-sharpness tranche: S1--S4}

The separate threshold checks have a completely specified enumeration:
for each $p=1,\ldots,18$, the $p+1$ values
\begin{equation}
 E_p=\{2^{p+1}-1-2^j:0\leq j\leq p\}
 \label{eq:ep}
\end{equation}
are all the integers with $p$ set bits in positions $0,\ldots,p$. They
exhaust the integers of popcount $p$ with $n<2^{p+1}$, rather than all
integers in that interval. A higher bit would already exceed the threshold.
The report records the following four checks:
\begin{enumerate}
\item[\textbf{S1.}] At $n=T_p$, for every $p=1,\ldots,18$, the candidate
union has $2^p$ distinct values, no collisions, and minimum $2^{p-1}$.
The symmetry count is $2^{p+1}$. The report also specifies a pinned
\texttt{math.comb} comparison for $n\leq1000$. This is a cutoff for that
comparison, not a declaration that every row up to $1000$ was simulated.
Among these threshold generations it includes $p\leq9$ ($T_9=767$),
and excludes $p\geq10$ ($T_{10}=1535$).
\item[\textbf{S2.}] Each complete list $E_p$ contains exactly two values
at most $T_p$, namely $2^p-1$ and $T_p$, and no value strictly between
them. The $p=18$ cross-consistency check is the 19-value list
\eqref{eq:e18}.
\item[\textbf{S3.}] At each of the 18 Mersenne generations $n=2^p-1$,
the union has $2^p$ distinct values, with zero produced once at $k=L(n)$.
The recorded symmetry count is $2^{p+1}-1$.
\item[\textbf{S4.}] At every $n\in E_p$, for every $p=1,\ldots,18$,
the minimum produced value is $n-2k_{\max}-1$, with $k_{\max}=L(n)$.
\end{enumerate}
The report states that an earlier expectation of additional failing
generations below $T_p$ was refuted by the enumeration: the proposed gap
has no $p$-bit integers. The finding concerns ordering and the candidate
union; it is not a new failure mode for the automaton.

In particular, the endpoint union properties were verified computationally
for $n=T_p$ and $n=2^p-1$, $p=1,\ldots,18$, as reported in S1 and S3.
This statement has those finite domains only.

The threshold tranche supplies no additional coverage assertion for
$p>18$ or for arbitrary generations above $T_p$. The elementary ordering
argument in \eqref{eq:second} is logically distinct from its finite
enumeration checks. Likewise, direct counting of $U(n)$ is distinct from
simulating Rule 126 or establishing the missing equality
$B_n=\widehat B_n$.

\section{Prior art and the limits of the novelty case}
\label{sec:prior}

The submask/Lucas mechanism is folklore-grade. Lucas's theorem supplies the
parity criterion, and the binary factorization of Pascal's triangle and
its relation to additive cellular automata are standard ingredients. This
draft claims no novelty for those ingredients.

There is also specific prior art beyond that general mechanism:
\begin{itemize}
\item A267442 already records the positive-index identity
\eqref{eq:headline} as Mathar's 2016 conjecture~\cite{oeis267442}.
Wolfram discusses the Rule 129 center-column pattern in 2002~\cite{wolfram129}.
\item The Rule 126 triangle A071035 already identifies A036987 as its
central column~\cite{oeis071035}.
\item A071051 already gives the Rule 126 row count
$2^{1+\pc(n)}$, with one subtracted at Mersenne generations, in a comment
by N.\ J.\ A.\ Sloane dated 9 August 2014~\cite{oeis071051}. Sloane also
states this formula in Section 6.1 of his 2015 paper~\cite{sloane}.
\item A071035 identifies the triangle A127872 as the right half of its
Rule 126 array~\cite{oeis071035}. A127872 includes G.\ C.\ Greubel's
18 April 2017 code for the formula
\[
 T(n,k)=\left(\binom{2n-1}{n-k}-\binom{2n-1}{n-k-2}\right)\bmod2
 \quad(0\leq k\leq n,\ n>0),
\]
with $T(0,0)=1$ and out-of-range binomial coefficients set to
zero~\cite{oeis127872}. This is prior art for a binomial representation,
even though the entry does not supply the needed dynamical derivation.
\end{itemize}
The last formula is already algebraically equivalent to the candidate
\eqref{eq:binomial} on the right half. Put $m=n-k$. Two applications of
Pascal's identity, followed by binary Lucas, give
\begin{equation}
 \binom{2n-1}{m}-\binom{2n-1}{m-2}
 \equiv\binom{2n+1}{m}
 \equiv\binom{n}{\lfloor m/2\rfloor}\pmod2.
 \label{eq:prior-equivalence}
\end{equation}
The first congruence uses that the middle Pascal coefficient is even and
that plus and minus agree modulo two. Thus the algebraic comparison is
available already; the remaining obligation is the dynamical derivation
and its formalization, rather than this rewriting.

The precise lower-part union formulation \eqref{eq:union} appears
unpublished in the limited sources inspected for this draft. This is a
provisional observation about presentation, not a priority claim: it is an
elementary reformulation of the catalogued binomial expression, conditional
on that expression's identification with the automaton. The identity
itself is already published as a conjecture, so the stronger wording that
both the characterization and identity are unpublished would be
inaccurate. A completed machine-checked derivation may be a useful
contribution, but mathematical novelty is not established by this draft or
by the finite evidence. The limited source check did not review all linked
manuscripts, including A.\ J.\ Macfarlane's manuscript linked from
A071051~\cite{oeis071051}.

\section{Retained draft follow-up requirements}

The following paragraph preserves the earlier draft's work list and its
then-pending status. The completed pinned proof is described in
Section~\ref{sec:formal-proof}; the requirements for raw computational
records and wider literature assessment remain unchanged.

The next mathematical deliverable is the missing recurrence-to-closed-form
argument, followed by a Lean statement and checked derivation with the
same seeds and indexing as this paper. The raw tranche reports and scripts
are needed to replace the missing entries in Table~\ref{tab:tiers} and to
document which computations independently check the dynamics. Further
literature review is needed before any novelty claim is finalized; the
equivalence \eqref{eq:prior-equivalence} already limits such a claim.
Until those items are supplied, this manuscript is an
evidence-based draft with a target theorem and a proof roadmap.

\begin{thebibliography}{99}
\raggedright

\bibitem{oeis267442}
OEIS Foundation Inc., \emph{The On-Line Encyclopedia of Integer Sequences},
entry A267442, central column of the single-seed Rule 129 evolution.
Includes R.\ J.\ Mathar's conjecture, 17 January 2016.
\url{https://oeis.org/A267442}.

\bibitem{oeis036987}
OEIS Foundation Inc., \emph{The On-Line Encyclopedia of Integer Sequences},
entry A036987, Fredholm--Rueppel sequence.
\url{https://oeis.org/A036987}.

\bibitem{oeis043545}
OEIS Foundation Inc., \emph{The On-Line Encyclopedia of Integer Sequences},
entry A043545, complement of A036987.
\url{https://oeis.org/A043545}.

\bibitem{oeis071035}
OEIS Foundation Inc., \emph{The On-Line Encyclopedia of Integer Sequences},
entry A071035, Rule 126 triangle; central-column cross-reference.
\url{https://oeis.org/A071035}.

\bibitem{oeis071051}
OEIS Foundation Inc., \emph{The On-Line Encyclopedia of Integer Sequences},
entry A071051, Rule 126 row counts; formula comment by N.\ J.\ A.\ Sloane,
9 August 2014.
\url{https://oeis.org/A071051}.

\bibitem{oeis127872}
OEIS Foundation Inc., \emph{The On-Line Encyclopedia of Integer Sequences},
entry A127872, right half of the Rule 126 triangle.
\url{https://oeis.org/A127872}.

\bibitem{sloane}
N.\ J.\ A.\ Sloane, \emph{On the Number of ON Cells in Cellular Automata},
2015, Section 6.1. arXiv:1503.01168v1.
\url{https://arxiv.org/abs/1503.01168}.

\bibitem{lucas}
\'Edouard Lucas, \emph{Sur les congruences des nombres eul\'eriens et des
coefficients diff\'erentiels des fonctions trigonom\'etriques suivant un
module premier}, Bulletin de la Soci\'et\'e Math\'ematique de France
\textbf{6} (1878), 49--54.
\href{https://doi.org/10.24033/bsmf.127}{doi:10.24033/bsmf.127}.
\url{https://www.numdam.org/item/BSMF_1878__6__49_1/}.

\bibitem{wolfram126}
Stephen Wolfram, \emph{A New Kind of Science}, Wolfram Media, 2002,
p.\ 55 (elementary rules) and pp.\ 269--270 (Rule 126 and Rule 90).
\url{https://www.wolframscience.com/nks/p55--more-cellular-automata/};
\url{https://www.wolframscience.com/nks/p269--special-initial-conditions/}.

\bibitem{wolfram129}
Stephen Wolfram, \emph{A New Kind of Science}, Wolfram Media, 2002,
p.\ 641, discussion of the center column of Rule 129.
\url{https://www.wolframscience.com/nks/p641--computations-in-cellular-automata/}.

\bibitem{report-consolidation}
\emph{Own-harness research---Rule 126 popcount series consolidation},
27 October 2026, orchestrator run approximately 19:53 PDT.
Unpublished supplied evidence summary, \texttt{REPORT-consolidation.md}.

\bibitem{report-threshold}
\emph{Own-harness tranche---Rule 126 popcount threshold sharpness},
27 October 2026, run approximately 20:23 PDT.
Unpublished supplied evidence summary, \texttt{REPORT-threshold.md}.

\end{thebibliography}

\noindent\small Online references consulted 28 October 2026. The reports
are local source documents, not public proof artifacts.

\end{document}
